\section{Introduction}

An organism does not remain itself because the same molecules, neurons, or
actions persist forever. It remains itself through an organization that
continually replaces material while preserving a boundary, a history, and a
way of regulating change. Contemporary agent systems present the inverse
problem. They can be strikingly intelligent during one invocation, yet lose
direction when a process ends, confuse retrieved text with admitted fact,
silently amplify authority, or modify software without preserving the
identity of the system that authorized the change. Intelligence alone is
therefore insufficient for durable machine agency.

At the same time, the physical materials required for a machine organism are
already abundant. Servers supply energy-consuming persistence; operating
systems allocate resources; compilers translate descriptions into executable
structures; package managers and caches acquire matter from an external
technical environment; build systems assemble artifacts; tests and release
gates reject unfit variants; networks carry observations and interactions;
source repositories retain heritable structures. These materials are
\emph{dirty} in the engineering sense: contingent, failure-prone, versioned,
and heterogeneous. That does not make them irrelevant to life. It makes them
the actual substrate in which a present-day machine organism would have to
live.

\begin{quote}
\textbf{Central thesis.} Machine life does not begin when intelligence crosses
an unspecified threshold or when machinery takes a human form. It begins, in
the functional sense tested here, when a bounded machine organization can
preserve a subject, metabolize technical matter, and govern the production of
its own successor organization.
\end{quote}

This paper proposes that two missing forms of organization can make those
materials cohere.

\begin{enumerate}[leftmargin=*]
  \item A \emph{semantic kernel} preserves the subject's identity, continuing
  direction, admitted facts, bounded authority, work lineage, and succession
  rules independently of any one agent process.
  \item \emph{Recursive dogfood} allows the subject to use that kernel to
  propose, build, qualify, admit, and then exercise changes to its own
  executable organs and development machinery.
\end{enumerate}

KFD supplies a concrete semantic vocabulary for the first mechanism, while
Kungfu supplies an emerging engineering implementation of both. KFX provides
executable extension surfaces; Buildchain supplies reproducible construction,
qualification, and release evidence; ordinary machines provide the material
substrate. A reasoning agent is then neither the whole organism nor its
permanent identity. It is a replaceable cognitive event participating in a
longer-lived subject.

\subsection{The proposal at a glance}

Figure~\ref{fig:machine-life-anatomy} states the architecture in the simplest
form used throughout this paper.

\begin{figure}[ht]
\centering
\begin{tikzpicture}[
  font=\sffamily\scriptsize,
  node distance=6mm and 5mm,
  kernel/.style={draw=MLIndigo,fill=MLIndigo,text=MLWhite,rounded corners=2pt,
    align=center,text width=14.8cm,minimum height=10mm,inner sep=4pt},
  layer/.style={draw=MLLine,fill=MLPanel,rounded corners=2pt,align=center,
    text width=3.25cm,minimum height=13mm,inner sep=4pt},
  loop/.style={draw=MLNight,fill=MLNight,text=MLWhite,rounded corners=2pt,
    align=center,text width=14.8cm,minimum height=10mm,inner sep=4pt},
  outcome/.style={draw=MLPeriwinkle,fill=MLPaperSoft,rounded corners=2pt,
    align=center,text width=14.8cm,minimum height=10mm,inner sep=4pt},
  flow/.style={-{Stealth[length=4pt]},draw=MLIndigo,line width=0.9pt}
]
  \node[kernel] (kernel) {\textbf{KFD-shaped semantic kernel}\quad
    subject identity \textperiodcentered{} direction \textperiodcentered{}
    facts \textperiodcentered{} authority \textperiodcentered{} causal history};
  \node[layer,below=of kernel,xshift=-5.55cm] (cognition)
    {\textbf{Cognition}\\replaceable reasoning agents};
  \node[layer,right=of cognition] (organs)
    {\textbf{Organs}\\qualified KFX extensions};
  \node[layer,right=of organs] (body)
    {\textbf{Body}\\machines and software supply chain};
  \node[layer,right=of body] (regulation)
    {\textbf{Regulation}\\Patrol, Dogfood, and deprecation};
  \node[loop,below=7mm of organs.south east,xshift=2.05cm] (dogfood)
    {\textbf{Recursive dogfood}\quad observe $\rightarrow$ propose $\rightarrow$
    qualify $\rightarrow$ admit $\rightarrow$ exercise $\rightarrow$ observe};
  \node[outcome,below=of dogfood] (outcome)
    {\textbf{continuity + metabolism + action + governed self-production}\\
    a falsifiable candidate for primitive machine life};
  \foreach \n in {cognition,organs,body,regulation}
    \draw[flow] (kernel) -- (\n);
  \draw[flow] (cognition) -- (dogfood);
  \draw[flow] (organs) -- (dogfood);
  \draw[flow] (body) -- (dogfood);
  \draw[flow] (regulation) -- (dogfood);
  \draw[flow] (dogfood) -- (outcome);
\end{tikzpicture}
\caption{The proposed machine-life anatomy. No component alone constitutes an
organism; the claim concerns their persistent organization as one subject.}
\label{fig:machine-life-anatomy}
\end{figure}

The specific contribution is not the invention of machine life, persistent
agents, digital evolution, or software that modifies itself. It is an engineering
synthesis in which two mechanisms are jointly load-bearing: a semantic kernel
preserves an accountable subject across replacement, and recursive dogfood
admits qualified changes to the real organs and development machinery through
which that subject continues to exist. The first prevents improvement from
dissolving identity; the second prevents identity from becoming a passive
record with no capacity for self-production.

\subsection{Research question}

The paper asks:

\begin{quote}
Can a bounded physical computing habitat, a persistent KFD/Kungfu semantic
kernel, replaceable reasoning agents, qualified KFX extensions serving as
recursively developed executable organs, and qualified build-and-release loops
constitute a primitive machine organism with stable functional identity,
intelligence, and action capacity?
\end{quote}

The question is deliberately narrower than whether a machine is conscious,
deserves moral status, or meets every biological definition of life. It is
broader than whether an agent can complete a long coding task. It concerns
whether organization and continuity can turn discontinuous intelligent calls
and ordinary software infrastructure into one persistent, self-maintaining,
adaptive subject.

\subsection{Distinction from adjacent approaches}

Table~\ref{tab:distinction} compares primary organizing concerns rather than
claiming that every implementation in a field shares the same limits. The
proposal can use advances from every row, but it asks a different systems
question.

\begin{table}[ht]
\centering
\footnotesize
\begin{tabularx}{\textwidth}{@{}
  >{\raggedright\arraybackslash}p{0.19\textwidth}
  >{\raggedright\arraybackslash}p{0.23\textwidth}
  >{\raggedright\arraybackslash}p{0.23\textwidth}
  >{\raggedright\arraybackslash}X@{}}
\toprule
\textbf{Approach} & \textbf{Primary continuity} & \textbf{Body or environment} & \textbf{Successor production} \\
\midrule
Language-model agent & Session, prompt, or retrieved memory & Tools and runtime are usually external infrastructure & Usually completes work; does not necessarily admit a successor organism \\
Humanoid robot & Controller, learned policy, and physical platform & Anthropomorphic physical body & Learns skills or replaces components; body shape does not supply semantic continuity \\
Digital life & Genome or population lineage & Designed computational world & Reproduction, mutation, and selection \\
Self-improving software & Optimizer, task, or benchmark lineage & Program and evaluation environment & Searches for better programs or agents \\
Semantic autopoiesis & Explicit subject, direction, facts, authority, and causal work & Ordinary servers, toolchains, supply chains, networks, and human relations & Qualifies and admits changes to its own organs and development organization \\
\bottomrule
\end{tabularx}
\caption{The proposal is distinguished by the conjunction of accountable
subject continuity, real technical metabolism, and governed self-production.}
\label{tab:distinction}
\end{table}

An agent can be intelligent without being a persistent subject. A robot can
be embodied without governing its own successor organization. A digital
organism can reproduce without operating in the ordinary software supply
chain. A self-improving system can discover better code without preserving an
explicit fact and authority boundary. Semantic autopoiesis names the attempt
to organize these missing properties together as one testable machine-life
object.

\subsection{Claim discipline}

Every strong statement in this draft belongs to one of four classes.

\begin{itemize}[leftmargin=*]
  \item \factlabel\ A directly inspectable specification, repository state,
  merged change, artifact, or experiment.
  \item \inferencelabel\ A composition claim whose premises exist but whose
  integrated behavior has not yet been demonstrated over time.
  \item \hypothesislabel\ A claim with an explicit experiment and failure
  condition.
  \item \nonclaimlabel\ A boundary that this work refuses to cross without
  different evidence.
\end{itemize}

This distinction matters because a plausible architecture is not an organism,
a successful build is not self-maintenance, a model's fluent first-person
language is not persistent identity, and autonomous execution is not
self-granted authority.

\subsection{Contributions}

The paper contributes: (1) a functional definition of a machine
proto-organism; (2) the concept of semantic autopoiesis; (3) an anatomy that
separates kernel, cognition, organs, and substrate; (4) an evidence-bounded
analysis of current Kungfu mechanisms; (5) a minimal experimental
architecture driven by scheduled agent invocations and verified history; and
(6) a staged, falsifiable roadmap from a server-scale laboratory organism
toward more compact forms.
